\documentclass[border=10pt]{standalone}
\usepackage{fontspec}
\AtBeginDocument{\hyphenation{zenodo}}% never hyphenate "zenodo" inside a DOI
\setmainfont{Noto Serif}[Ligatures=TeXOff]
\XeTeXgenerateactualtext=1
\newfontfamily\bengalifont{Noto Serif Bengali}[Script=Bengali, NFSSFamily=bnfam]
\usepackage[Latin,Bengali]{ucharclasses}
% group-free font switching (as in header-common.tex): robust inside braces, cells and quotes
\makeatletter
\def\BN@saved{\rmdefault}
\def\BN@bn{bnfam}
\newcommand\BNin{\edef\BN@tmp{\f@family}\ifx\BN@tmp\BN@bn\else\let\BN@saved\BN@tmp\fi\fontfamily{bnfam}\selectfont}
\newcommand\BNout{\fontfamily{\BN@saved}\selectfont}
\makeatother
\setTransitionsFor{Bengali}{\BNin}{\BNout}
% Typographic quotes and dashes are not in a ucharclasses class (class 0), so a Bangla word right
% after “ or ‘ would not trigger the Bangla font. Treat them as Latin so that every switch fires.
\XeTeXcharclass"2018=\XeTeXcharclass`A
\XeTeXcharclass"2019=\XeTeXcharclass`A
\XeTeXcharclass"201C=\XeTeXcharclass`A
\XeTeXcharclass"201D=\XeTeXcharclass`A
\XeTeXcharclass"2013=\XeTeXcharclass`A
\XeTeXcharclass"2014=\XeTeXcharclass`A
\XeTeXcharclass"2026=\XeTeXcharclass`A
\usepackage{newunicodechar}
\newunicodechar{।}{{\bengalifont\char"0964}}
\newunicodechar{॥}{{\bengalifont\char"0965}}

\usepackage[table]{xcolor}
\usepackage{booktabs,tabularx,array}
\definecolor{bapilrule}{HTML}{1F2937}
\definecolor{bapilstripe}{HTML}{F3F4F6}
\renewcommand{\arraystretch}{1.28}
\setlength{\tabcolsep}{5pt}

\begin{document}
\begin{minipage}{17cm}
{\large\bfseries সারণি ৩৯।}\ {\large ঝুঁকি-স্তর এবং প্রতিটি স্তরের ন্যূনতম আবশ্যিক নিয়ন্ত্রণ}\par\vspace{7pt}
\small
\rowcolors{2}{bapilstripe}{white}
\begin{tabularx}{\linewidth}{@{}>{\hsize=0.667\hsize\raggedright\arraybackslash}X>{\hsize=1.067\hsize\raggedright\arraybackslash}X>{\hsize=1.600\hsize\raggedright\arraybackslash}X>{\hsize=0.667\hsize\raggedright\arraybackslash}X@{}}
\arrayrulecolor{bapilrule}\toprule
\rowcolor{white}\textbf{স্তর} & \textbf{সাধারণ ব্যবস্থা} & \textbf{ন্যূনতম আবশ্যিক নিয়ন্ত্রণ} & \textbf{অনুমোদনকারী} \\
\midrule
\textbf{T0} নিষিদ্ধ & সামাজিক স্কোরিং; গণ বা নির্বিচার বায়োমেট্রিক নজরদারি; মানুষের রাজনৈতিক বা ধর্মীয় মতের প্রোফাইলিং; বাস্তব ব্যক্তির প্রতারণামূলক কৃত্রিম মিডিয়া; দুর্বলতাকে কাজে লাগিয়ে প্রভাবিত করা & কোনো ব্যবস্থাতেই BAPIL এগুলো নির্মাণ, পাইলট বা সমর্থন করে না & সনদ দ্বারা বর্জিত \\
\textbf{T1} নিম্ন & উন্মুক্ত নথিতে গবেষণা ও অনুসন্ধান সরঞ্জাম; প্রকাশিত ডেটার অভ্যন্তরীণ ড্যাশবোর্ড & এআই ব্যবস্থার রেকর্ড; মডেল বা ডেটা কার্ড; মূল্যায়ন প্রতিবেদন; নিরীক্ষা লগ; স্বচ্ছতা নিবন্ধনে প্রকাশ্য তালিকাভুক্তি & ইউনিট প্রধান ও TAB \\
\textbf{T2} সীমিত & কর্মকর্তাদের ব্যবহৃত পরামর্শমূলক সরঞ্জাম; প্রকাশ্য ড্যাশবোর্ড, প্রতিবেদন ও অন্যান্য জনমুখী তথ্য সরঞ্জাম; পরিচয়-অপসারিত মতামতের সামষ্টিক বিশ্লেষণ & T1-এর সব নিয়ন্ত্রণ এবং: ERRC পর্যালোচনা; রেড-টিম পরীক্ষা; ব্যবহারে প্রতিটি ফলাফলের মানব-পর্যালোচনা (প্রকাশ্য ড্যাশবোর্ড ও প্রতিবেদনের ক্ষেত্রে প্রতিটি সংস্করণের পর্যালোচনা); ব্যবহারকারী প্রশিক্ষণ; ঘটনা-প্রক্রিয়া; ত্রৈমাসিক পর্যবেক্ষণ প্রতিবেদন & ERRC (স্টিয়ারিং বোর্ডের পরামর্শসহ) \\
\textbf{T3} উচ্চ & নির্দিষ্ট ব্যক্তি বা প্রতিষ্ঠান-সংক্রান্ত সিদ্ধান্তে ভূমিকা রাখে, বৃহৎ পরিসরে ব্যক্তিগত ডেটা প্রক্রিয়া করে বা প্রয়োগকারী কাজে ব্যবহৃত যেকোনো ব্যবস্থা & BAPIL T3 পরিচালনাগত ব্যবস্থা নির্মাণ করে না। শুধু গবেষণা, এই শর্তে: DPIA, আইনি পর্যালোচনা, স্বাধীন নিরীক্ষা, ERRC অনুমোদন এবং প্রকাশিত প্রভাব মূল্যায়ন & ERRC (বাধ্যতামূলক) ও স্বাগতিক সংস্থার আইনি পর্যালোচনা \\
\bottomrule
\end{tabularx}
\par\vspace{3pt}{\scriptsize\color{gray} উৎস: BAPIL টেকনিক্যাল ব্লুপ্রিন্ট, সংস্করণ ১.০ (২০২৬)। CC BY 4.0।}
\end{minipage}
\end{document}
